\section{The refutation special case: four first-order examples}
\label{sec:examples}

We now take the constant target $G_n=1$. This is the refutation special
case of the uniform-consequence theorem; the sharper invariant-target
budget applies. The four cases all occur among fixed first-order sentences with no finite
nonempty models. We give examples over a field of fixed prime characteristic
$p$. Each sentence is translated by the graph-and-sum convention: relation
and function-graph variables are Boolean, and every function has a totality
sum and pairwise functionality equations. The parameter $n\geq1$ is the
universe size. For the matrix notation in this section, we relabel the
universe as $\{1,\ldots,n\}$; this changes no equation or degree.

\paragraph{T1: a local contradiction.}
The sentence $\forall x\,(R(x)\wedge\neg R(x))$ contributes $r_i$ and
$1-r_i$. Their sum refutes the system in degree one, independently of $n$.
The contradiction is already visible at one element.

\paragraph{T2: an injection that misses a point.}
Let $\eta$ say that a unary function $g$ is injective and that a constant
$c$ is outside its image. Write $x_{ij}$ for $g(i)=j$ and $c_j$ for $c=j$.
In addition to the Boolean equations, its system $\Gamma_n^\eta$ consists of
\begin{equation}\label{eq:G-witness}
\begin{gathered}
 \sum_jx_{ij}-1=0,\qquad
 x_{ij}x_{ik}=0\ (j\ne k),\qquad
 x_{ij}x_{hj}=0\ (i\ne h),\\
 \sum_jc_j-1=0,\qquad
 c_jc_k=0\ (j\ne k),\qquad x_{ij}c_j=0.
\end{gathered}
\end{equation}
A solution would inject $n$ elements into the $n-1$ elements other than $c$,
so every instance is unsatisfiable. For either proof system, every
degree-$d$ refutation satisfies
\[
 n<3pd\qquad(n\geq3).
\]
Here is how the internal matching bound gives this estimate. Choose the
largest $p$-power $q\leq n/3$ and put $r=n-q$. Set $c_n=1$, all other
$c_j=0$, and restrict the image of $g$ to $[r]$. After transposition, every
remaining axiom is an axiom of the bijection system $B_{r,n}$.
Here $B_{r,m}$ is the Boolean $r$-by-$m$ matrix system with every row and
column sum equal to one and both row and column collisions equal to zero.
If $d<q$, then $2d\leq r$ and $q\mid n-r$, contradicting the rectangle
obstruction of Lemma~\ref{lem:rectangle-obstruction}. The substitutions do
not increase ordinary-PC or NS degree. Thus $d\geq q$, while maximality of
$q$ gives $n<3pq\leq3pd$. This proves T2 without an external pigeonhole
lower-bound theorem.

\paragraph{A matching linear upper bound.}
For the injection system, put
\[
 R_i=\sum_jx_{ij},\qquad C=\sum_jc_j,\qquad P=\prod_{i=1}^nR_i.
\]
Every monomial $c_j\prod_i x_{i,\lambda(i)}$ in the expansion of $CP$
is divisible by an input equation from \eqref{eq:G-witness}. If
$\lambda$ is not injective, it contains a column collision. If
$\lambda$ is injective, it is a permutation of the finite universe;
for the row $i$ with $\lambda(i)=j$, the monomial contains the omission
factor $c_jx_{ij}$. Thus $CP$ lies in the raw NS span at degree at most
$n+1$. The telescoping identity
\[
 CP-1=(C-1)+C\sum_{i=1}^n(R_i-1)\prod_{h<i}R_h
\]
expresses $CP-1$ using the constant and row totality equations, with
every product of degree at most $n+1$. Subtraction gives a refutation,
so for every $n\ge1$,
\[
 d_{\PC}(\Gamma_n^\eta)\le d_{\NS}(\Gamma_n^\eta)\le n+1.
\]
This argument works over any field and uses no Boolean simplification.
Together with the preceding lower bound, it shows that both degrees are
$\Theta(n)$ for each fixed characteristic. The general linear gap is
therefore sharp even among all-size inconsistent first-order families.

\paragraph{T3: hardness on one residue class.}
Let $\chi_p$ assert
\[
 \forall x\qquad f^p(x)=x\quad\hbox{and}\quad f(x)\ne x.
\]
Its finite models are disjoint unions of directed $p$-cycles. In particular,
it has a $0/1$ graph assignment exactly when $p\mid n$. An explicit
normalisation uses the function equations for $f$, the equations $f_{ii}=0$,
and
\begin{equation}\label{eq:cycle-path}
 f_{i_0i_1}f_{i_1i_2}\cdots f_{i_{p-1}i_p}=0
 \qquad(i_p\ne i_0).
\end{equation}
If $p\nmid n$, summing all length-$p$ walks gives an NS refutation of degree
at most $p$. Indeed, totality derives the sum of their monomials minus $n$
within degree $p$. Nonclosed walks vanish by \eqref{eq:cycle-path}.
A closed walk with a repeated vertex either contains a functionality
collision or follows a shorter deterministic cycle; since $p$ is prime,
the latter is a forbidden fixed point. Each remaining cycle contributes its
monomial $p$ times, according to its starting vertex, and hence contributes
zero. The resulting nonzero scalar $n$ yields the refutation.

Now take $\psi_3=\eta\wedge\chi_p$ in disjoint nonlogical symbols.
For $p\mid n$, substitute a $p$-cycle assignment into the second block.
This preserves proof degree and shows that the conjunction has exactly the
degree of $\eta$, separately for NS and PC. For $p\nmid n$, the second
block refutes in degree at most $p$. Consequently this all-size
unsatisfiable family is bounded off the multiples of $p$ and has linear
degree on the multiples: it has type T3. No exact degree is needed on the
bounded residue classes.

\paragraph{T4: the valuation of the universe size controls the difficulty.}
For a positive integer $n$, its $p$-adic valuation $v_p(n)$ is the
largest $a$ such that $p^a$ divides $n$. Thus $n=p^a u$ with
$p\nmid u$. We use this integer statistic, not a valuation on the
coefficient field; the corresponding block size is $p^a$.
The weak-bijection example from the earlier classification~\cite[Section~5]{Riis}
is more revealing; its algebraic complexity is closely related to the
bijective pigeonhole principle studied by Beame and Riis~\cite{BeameRiis}.
Let unary functions $f,g$
be injective with disjoint images, and require their images together to
cover the universe. Use a unary witness $h$ for coverage:
\[
 \forall z\qquad f(h(z))=z\ \vee\ g(h(z))=z.
\]
Let $a_{ij},b_{ij},w_{ji}$ represent $f(i)=j$, $g(i)=j$, and $h(j)=i$.
The combined $2n$-by-$n$ matrix $(a;b)$ has row sums equal to one and both
row and column collisions equal to zero. The witness block consists of
\begin{equation}\label{eq:bhat-witness}
 \sum_iw_{ji}-1=0,\qquad
 w_{ji}w_{ji'}=0\ (i\ne i'),\qquad
 w_{ji}(1-a_{ij})(1-b_{ij})=0,
\end{equation}
together with the Boolean equations. Call this complete system
$\widehat B_n$. Its witnesses are retained: it is not identified with the
simpler matrix system $B_{n,2n}$.

Put $q=p^{v_p(n)}$. The proved bounds give
\begin{equation}\label{eq:bhat-window}
 q\leq d_{\mathrm{PC}}(\widehat B_n)
 \leq d_{\mathrm{NS}}(\widehat B_n)\leq q+4
 \qquad\hbox{whenever } n\geq2q.
\end{equation}
The upper bound holds at every positive size, without the displayed size
condition. We give both degree-preserving reduction and certificate
explicitly.

For the lower bound transpose the combined matrix and substitute
$w_{ji}=a_{ij}+b_{ij}$ into every witness variable. Its totality becomes
the corresponding totality of $B_{n,2n}$. The identities
\begin{align*}
(a+b)^2-(a+b)&=(a^2-a)+(b^2-b)+2ab,\\
(a+b)(1-a)(1-b)&=-(1-b)(a^2-a)-(1-a)(b^2-b)
\end{align*}
put the substituted Boolean and coverage equations in the raw NS spans
at degrees two and three respectively. A substituted witness collision
expands into four column collisions of the raw combined matrix.
Thus every substituted input has a derivation within its own degree.
Substitution is affine: a multiplication by an old variable becomes a
sum of permitted multiplications by new variables, with the preceding
line still of degree at most $d-1$. A degree-$d$ ordinary-PC refutation
would therefore refute $B_{n,2n}$ within the same degree. The rectangle
obstruction excludes this when $2d\le n$ and $d<q$, yielding the lower
bound in \eqref{eq:bhat-window}.

For the upper bound fix a column $j$ and abbreviate
\[
 u_i=a_{ij}+b_{ij},\qquad T=\sum_i u_i,\qquad
 Z=\sum_i w_{ji}u_i.
\]
The coverage axiom and disjointness collision give the exact identity
\[
 Z-1=(\sum_iw_{ji}-1)
 -\sum_iw_{ji}(1-a_{ij})(1-b_{ij})
 +\sum_iw_{ji}a_{ij}b_{ij}.
\]
Consequently $Z-1$ lies in the degree-three NS span. If $x$ is one entry
of this column, then
\[
 Tx-x=(x^2-x)+\sum_{y\ne x}yx
\]
lies in the degree-two span, by Booleanity and column collisions.
Adding this for $a_{ij}$ and $b_{ij}$ gives $Tu_i-u_i$ in that span.
Finally
\[
 T-1=\sum_iw_{ji}(Tu_i-u_i)-(T-1)(Z-1)
\]
puts the missing column totality in the degree-four span. This retains
the coverage witnesses throughout the derivation.

For completeness, consider any matrix with row sums $R_i$ and column
sums $C_j$, with their differences from one in the NS span at degree
$e$, and with both collision families in degree two. For a fixed
$k$-element row set $S$, telescoping the product of its row sums gives
$\prod_{i\in S}R_i-1$ in degree at most $e+k$. Summing over $S$ gives
\[
 \sum_{|S|=k}\prod_{i\in S}R_i-\binom{\#\mathrm{rows}}k
\]
in that same span. Expand the products as partial functions from $k$
rows to columns. A noninjective partial function contains a column
collision, so its monomial is already in the degree-$k$ span.
The remaining monomials are exactly those of $k$-edge matchings.
The column calculation gives the same matching sum, since inversion
pairs partial injections with identical monomials. Subtracting derives
\[
 \binom{\#\mathrm{rows}}k-\binom{\#\mathrm{columns}}k
\]
within degree $e+k$. Applied to the $2n$-by-$n$ matrix with $e=4$, this
is the desired counting certificate.

Write $n=qu$ with $q=p^{v_p(n)}$ and $p\nmid u$. In characteristic
$p$, $(1+t)^n=(1+t^q)^u$. The coefficient of $t^q$ is therefore $u$;
for $(1+t)^{2n}$ it is $2u$. Thus
$\binom{2n}q-\binom nq=u\ne0$ in $K$. Rescaling the counting
certificate gives an NS refutation of degree at most $q+4$.

For each fixed $a$, the sizes $n=p^a(p(k+1)+1)$, $k\geq0$, have
$v_p(n)=a$ and satisfy \eqref{eq:bhat-window}. Hence some finite degree in
$[p^a,p^a+4]$ recurs infinitely often. These windows are unbounded as $a$
varies, proving T4 for both proof systems. On the other hand, $p\nmid n$
implies degree at most five. There are explicit simultaneous linear peaks:
at $n_a=(p+1)p^a$,
\[
 \frac{n_a}{p+1}\leq d_{\mathrm{PC}}(\widehat B_{n_a})
 \leq d_{\mathrm{NS}}(\widehat B_{n_a})
 \leq\frac{n_a}{p+1}+4.
\]
For example, in characteristic two every odd $n$ has degree at most
five, whereas at $n=3\cdot2^a$ both degrees lie between $2^a$ and
$2^a+4$. The peaks are linear in their universe size even though the
profile returns to bounded values infinitely often.
Thus linear peaks coexist with infinitely many easy sizes. This is the
reason the T4 conclusion concerns peaks rather than an eventual linear
lower bound at every size.
